\cotefile{70}
\scanpage{17}
% archive sheet number: 16
% source lines 1-9: block cancelled by long diagonal strokes
\[
\deleted{\sigma_i(x_1,\ldots,x_{n+1})=(x_1,\ldots,x_{i-1},\underset{!}{\underline{x_{i+1}}},x_i,\lambda_{i,i+2}x_{i+2},\lambda_{i,i+2}x_{i+3},\ldots,\lambda_{i,i+2}x_{n+1})}
\]
\[
\deleted{H_i:\left\{\begin{aligned}
x_i&=x_{i+1}\\
(\lambda_{i,i+2}-1)x_{i+2}&=\cdots=(\lambda_{i,i+2}-1)x_{n+1}=0
\end{aligned}\right.}
\]
\deleted{Si $H'_i$ est l'\unsure{hyperplan} d'équation $x_i=x_{i+1}$, on a $H_i\subset H'_i$, donc $H_i$ n'est un \unsure{hyperplan} que si $H_i=H'_i$, i.e. les équations $(\lambda_{i,i+2}-1)x_{i+2}=\cdots=(\lambda_{i,i+2}-1)x_{n+1}=0$ sont \unsure{conséquence|conséquences} de $x_i=x_{i+1}$, ce qui bien sûr n'est le cas que si $\lambda_{i,\unsure{i+1|i+2}}=1$.}

% source lines 10-13
\[
\left\{\begin{aligned}
&\lambda_{i,1}(x_{i+2}+\cdots+x_{n+1})=\lambda_{i,2}(x_{i+2}+\cdots+x_{n+1})=\cdots\ \lambda_{i,i-1}(x_{i+2}+\cdots+x_{n+1})=0\\
&\deleted{\unsure{x_{i+1}}}\lambda_{i,i}(x_{i+2}+\cdots+x_{n+1})=x_i-x_{i+1}\\
&\lambda_{i,i+1}(x_{i+2}+\cdots+x_{n+1})=x_{i+1}-x_i\\
&(\lambda_{i,i+2}-1)x_{i+2}=\cdots=(\lambda_{i,i+2}-1)x_{n+1}=0
\end{aligned}\right.\tag*{$(H_i)$}
\]

% source lines 14-16
(en plus, les conditions sur les $\lambda_{i,j}$, $\unsure{0|1}\leq j\leq i+2$, \deleted{\unsure{disant que}} \unsure{exprimant} $\sigma_i^2=\mathrm{id}\ \ldots$)

% source lines 17-19
\deleted{\unsure{Si}} Sur un corps: si on avait $\lambda_{i,i+2}\ne1$, comme $\lambda_{i,i+2}^2=1$, on aurait $\lambda_{i,i+2}=-1\ne1$ (donc car $\ne2$), et les équations de $H_i$ deviennent
\[
\left\{\begin{aligned}
&x_{i+2}=\cdots=x_{n+1}=0\\
&\deleted{\lambda_{i,1}\unsure{\overline{\Sigma}_i|\Sigma_i}(x)=\cdots=\lambda_{i,i-1}\overline{\Sigma}_i(x)=0}\quad\deleted{:\ \overline{\Sigma}_i(x)=x_{i+2}+\cdots+x_{n+1}}\\
&\deleted{\lambda_{i,i}\overline{\Sigma}_i(x)=}\ x_i=x_{i+1}\ \unsure{\}}\\
&\deleted{\lambda_{i,i+1}\overline{\Sigma}_i(x)=x_{i+1}-x_i}
\end{aligned}\right.
\]
% source lines 24-29
ce qui fait $[n+1-(i+1)]+1=n-i+1$ équations \deleted{\unsure{Si} $i\leq n+1$, les premières équations \unsure{étant} $n+1-i$} indépendantes et comme $i\leq n-1$, on trouve une sous-variété linéaire de codim $\geq2$.

Donc on doit avoir $\lambda_{i,i+2}=1$, les équations de $H_i$ deviennent
